% ══════════════════════════════════════════════════════════════════
% 第一章 绪论
% ══════════════════════════════════════════════════════════════════
\chapter{绪\hspace{6pt}论}

\section{研究背景与意义}

大型工业基础设施——桥梁、压力容器、风力发电机组和石油天然气管道——在其服役周期内持续承受疲劳载荷、腐蚀环境和意外冲击，表面及近表面裂纹、壁厚减薄和材料退化是其最常见的损伤模式。对这类结构实施定期无损检测并非可选维护策略，而是保证公共安全和避免灾难性失效的强制性工程需求\citing{qu2020development}。然而，受检结构动辄数十米高的作业面、受限空间内不可达的检测点、以及高空作业固有的安全风险，使传统人工搭架检测在效率、成本和人员安全三个维度上均面临难以调和的矛盾\citing{ball1998detection}。

多旋翼无人机为上述困境提供了独特的解决方案。与固定翼或大型直升机平台不同，多旋翼的悬停能力、低空机动性和对复杂几何表面的适应性使其成为近距离结构检测的理想空中载体。过去十年间，无人机平台从简单的航拍摄像发展为搭载高清可见光相机、红外热像仪和激光雷达的多传感器系统，实现了从远距离目视筛查向近场定量测量的功能跃迁\citing{dou2025uav}。然而，视觉与热成像等非接触传感模态共享一个根本局限：它们只能探测表面或近表面（对热波而言）异常，对隐藏在材料内部的亚表面分层、壁厚减薄和体缺陷缺乏穿透能力。

电磁超声换能器（EMAT）的引入打破了上述深度限制。EMAT利用交变线圈在导电材料集肤深度内感生涡流，涡流与永磁体提供的静偏置磁场耦合产生洛伦兹体力，后者直接作为体载荷源在材料内部激发超声波\citing{sun2025emat}。这一换能机理带来了两个关键优势：第一，超声波在工件内部而非探头中产生，消除了传统压电超声对液体耦合剂的依赖，使检测可在干燥、高温和粗糙表面上进行；第二，回波信号携带了完整的传播路径信息——从激励线圈的涡流分布、到材料内部的弹性波传播、再到底部界面的反射——由此可同时提取壁厚、缺陷位置和材料声学属性。在静态实验室条件下，基于EMAT的测厚精度已达到亚毫米级，测量误差低于1\%\citing{kubrusly2023measurement}，且在500°C高温和5 mm提离距离下仍能保持优于3.5\%的精度\citing{chen2024novel}。

然而，EMAT的上述精度承诺建立在探头与被测表面之间保持稳定相对位置的前提之上——这一前提在无人机搭载条件下恰恰被根本性地打破了。问题的核心在于：EMAT的换能效率对探头-表面间距（升距）极其敏感。升距的微小变化会通过两种机制同时恶化信号质量：磁通密度随距离按负幂律衰减，导致涡流密度和洛伦兹力幅值下降；涡流的空间分布也因磁力线扩散而改变，使激励声场的指向性和模态纯度退化。当EMAT探头被刚性安装在多旋翼机体上时，平台自身的姿态微振（±2°量级）、螺旋桨尾流引起的不稳定地面效应、以及接触瞬间的动力学反冲，都会在亚毫米尺度上持续调制升距——而这一调制尺度恰好与EMAT的敏感区间重叠。

更根本的困难在于，上述升距变化并非以孤立方式发生。无人机的法向速度收敛过程、末端执行器的力学刚度、被测表面的曲率和材料属性，以及环境风场的随机扰动，共同构成了一个强耦合的多因素系统。在该系统中，接触状态的改变不是简单的二值事件，而是升距逐渐减小、法向速度渐近收敛、回波能量非线性增长和频谱特征持续演化的连续动态过程。单一传感器——无论是超声回波幅值、飞控位姿残差还是视觉深度估计——都无法独立可靠地表征这一复杂过程。当升距因姿态扰动而瞬时增大时，回波幅值的下降既可能源于真实的接触脱离，也可能来自暂时性的磁耦合减弱，而后者在姿态恢复后会自动消失。依赖单帧幅值阈值的传统判别方法在物理上无法区分这两种场景，其误判不可避免地导致控制闭环中的频繁状态切换——即接触-脱离-再接触的循环振荡。

上述分析指向一个清晰的研究空白：无人机EMAT检测需要一个能够联合解释电磁-弹性耦合、飞行器刚体动力学和视觉几何约束的多模态接触感知框架，而目前的方法论——无论是纯数据驱动的黑箱分类器，还是基于单一物理量阈值的工程判据——均不具备在动态扰动下维护物理一致性的能力。纯数据驱动方法缺乏对EMAT换能链中物理因果关系的显式建模，在传感器噪声增大或环境条件改变时泛化能力急剧下降；物理约束方法（如基于PINN的动力学建模）虽然引入了方程正则化，但大多局限于单一连续系统的状态估计，未涉及跨模态、跨时间尺度的异步观测融合问题。

本文正是在这一研究空白处开展工作。我们提出一种基于多模态逐级物理约束的框架，将EMAT超声回波序列、无人机六自由度位姿残差和RGBD局部视觉特征统一建模为一个时序感知系统。在该框架中，物理先验——包括时间邻近性、空间邻近性和法向一致性——被显式编码为Transformer注意力机制的对数偏置项，而非被隐含在数据分布中。这一设计使模型能够在飞行扰动下维持可解释的证据链：视觉确认几何可达性，位姿验证动力学收敛性，超声评估电磁耦合质量，三者同时满足时接触判定才被确立。该框架的理论价值在于将接触状态从工程经验判据提升为具有物理可证伪性的概率估计问题；其工程价值在于为无人机自主NDT提供了一条从传感器原始输出到闭环控制决策的完整推理路径，有望推动多旋翼平台从视觉巡检工具向真正的空中物理交互系统的范式转变。

\section{国内外研究现状}

\subsection{无人机搭载无损检测技术演进与核心瓶颈}

结构健康监测领域对空中检测平台的需求源于一个无法调和的工程矛盾：大型工业设施——桥梁缆索、风电叶片、储罐壁板和压力容器——的检测覆盖率要求与人工可达性之间存在根本性鸿沟。搭设脚手架或动用高空作业车不仅成本高昂（单次检测动辄数十万美元），而且在密闭空间、高耸结构和离岸设施等场景中存在致命安全风险\citing{qu2020development}。这一矛盾驱动了将多旋翼无人机用作移动传感平台的研究方向，其发展轨迹呈现出一条清晰的技术路线演化：从被动的远距离目视记录，到主动的物理激励与信号采集，再到自主决策与闭环操控。

\subsubsection{从视觉巡检到定量无损检测的范式转变}

无人机用于结构检测的早期实践——约2010至2015年——本质上是对人工目视检查的替代和延伸。搭载高清可见光相机的四旋翼平台使检测人员能够在地面安全位置获取结构表面的高分辨率图像，而红外热像仪则拓展了对近表面脱粘、水分侵入和隔热层缺陷的探测能力。这一阶段的系统架构相对简单：无人机作为可遥控的相机云台，数据采集与缺陷判别完全由地面操作人员完成。尽管该模式大幅提升了检测效率和人员安全性，但其功能边界也十分明确——成像传感器仅能捕获表面或热扩散可达的浅表层异常，对材料内部的分层扩展、壁厚减薄和体缺陷缺乏穿透性检测能力\citing{ball1998detection}。

为实现从表面筛查到内部缺陷定量的跃迁，研究者将目光投向了超声波、电磁涡流和激光超声等主动激励式无损检测模态。这一转变带来了全新的工程挑战：不同于视觉传感器可在数米乃至数十米的安全距离外工作，超声探头和涡流线圈要求与被测表面建立稳定的近场或直接物理接触。压电超声换能器需要耦合剂层实现声阻抗匹配，涡流检测对提离高度波动极其敏感，而激光超声则受限于激发效率和设备体积。无人机必须从“远端观察者”转变为“近端物理交互者”——这一角色转变涉及气动稳定性、接触力控制和末端定位精度等多个维度的系统性难题。

\subsubsection{接触式超声检测：从被动柔顺到主动力控}

接触式超声检测的空中实施面临的首要困难来自接触瞬间的动力学冲击。多旋翼无人机在自由飞行阶段依靠推力矢量维持姿态和位置，当末端探头突然与被测表面发生物理接触时，接触反力会瞬间打破力-力矩平衡——这一瞬态扰动的幅值可达数牛顿，持续时间数十毫秒\citing{gonzalez2019payload}。若不加专门处理，飞控系统会将接触力视为外部位置扰动并试图通过增大推力来补偿，反而导致探头与表面之间的接触力失控增长，造成检测失败甚至平台失稳。

应对这一问题的第一条技术路线是被动柔顺机构设计。通过在探头与机体之间串联弹簧-阻尼元件，接触冲击的动能被机械结构吸收和耗散，显著降低了传递至飞控回路的扰动幅值。此类方案结构简单、响应速度快（被动机械元件的响应时间在毫秒量级），适用于平面或小曲率表面上的初始接触建立。

然而，被动柔顺的局限性同样突出。弹簧的刚度-行程乘积决定了可吸收能量的上限，而这一乘积受探头质量和可用安装空间的双重约束；此外，弹簧恢复力会在接触维持阶段持续施加一个与形变成正比的推力，使得真正的法向接触力无法被独立控制。当检测要求恒定的法向接触力以维持稳定的声耦合条件时——例如压电探头需要的10-20 N等效压紧力——被动柔顺便力不从心。

第二条技术路线——主动力控——通过将接触力纳入无人机控制回路的反馈输入，实现了对法向力的闭环调节。Kocer等人\citing{kocer2019inspection}利用非线性模型预测控制将未知接触扰动作为外部力进行在线估计和补偿，在5.5 N恒定接触力下成功获取了天花板结构的超声A扫描数据。主动力控方案的关键优势在于其对不同表面刚度和几何形状的自适应能力：控制器可以根据力传感器的反馈动态调整推力输出，无需在检测前获知准确的表面机械阻抗。

\subsubsection{全驱动平台的推力矢量解耦与局限}

标准四旋翼属于欠驱动系统——其四个旋翼推力仅能独立控制垂直方向的升力和绕体轴的三个力矩，水平运动必须通过机体倾斜产生的推力分量实现。这一约束在接触式检测中产生了根本性的运动学冲突：当探头需要向垂直壁面施加法向力时，机体必须向壁面方向倾斜，而这一倾斜恰与保持探头法向对准壁面的控制目标相悖。换言之，欠驱动平台在接触检测中面临着推力方向与姿态指向之间的固有耦合，该耦合在曲面结构和倾斜表面上尤为突出。

全驱动和过驱动多旋翼平台通过引入可倾转旋翼或额外的独立推力执行器，将推力矢量的维数从4扩展至6，实现了位置与姿态的动力学解耦。Watson等人\citing{watson2022dry}利用Voliro全矢量推力无人机，在倾斜和悬垂表面上完成了干耦合超声轮式探头的连续B扫描，法向接触力维持在20 N左右。Marcellini等人\citing{marcellini2024development}在倾转旋翼平台上集成了双重阻抗控制——机械臂端负责精细力控，机体端负责姿态稳定——在接触阶段动态吸收物理冲击扰动。

然而，全驱动方案带来了新的代价：执行器数量的增加意味着更高的机械复杂度、更大的自重和更短的航时，而可倾转机构的旋转关节在工业现场的多尘和腐蚀性环境中面临可靠性挑战。更重要的是，上述研究——无论是被动柔顺、主动力控还是全驱动——均聚焦于从机械运动学角度实现并维持“物理接触”，而几乎未触及一个更深层的问题：物理接触的建立是否等同于有效的超声信号耦合？

\subsubsection{物理接触与信号耦合之间的鸿沟}

这一区分是理解无人机EMAT检测本质困难的关键。当压电超声探头通过干耦合弹性体与金属表面接触时，只要接触正压力超过一定阈值（通常为10-20 N），弹性体即可有效填充微观气隙，实现可重复的声阻抗匹配。然而，电磁超声换能器的耦合机制完全不同：EMAT的换能效率取决于探头线圈与被测表面之间的电磁场相互作用，而这一作用对线圈-表面间距——即升距——具有指数级的敏感性。升距每增加1 mm，涡流密度可下降30-50\%，而涡流密度的下降直接导致洛伦兹体力幅值的等比例衰减。

这一物理差异带来了一项无人机平台特有的工程困境。在压电检测中，接触力是声耦合的直接指标——力越大，耦合越充分，两者的关系单调且在相当宽的力区间内保持近线性。而在EMAT检测中，接触力本身并非耦合质量的可靠指标：即使法向力稳定，微小的姿态角扰动（±1°量级）也会通过改变线圈与表面的相对倾斜角而影响涡流的空间分布和磁力线的穿透深度。换言之，压电检测只需要维持一个标量（接触力），而EMAT检测需要同时维持一个包括升距、法向角和切向滑移在内的三维相对位姿约束。前者可由单输入单输出的力控环处理，后者则需要以多模态感知为基础的联合状态估计。

\subsubsection{多模态感知与物理约束融合的缺失}

无人机EMAT检测的上述独特性将当前研究的一个关键空白推至前台：飞行控制的反馈回路与检测信号的反馈回路本质上是异步、异质且缺乏统一解释框架的。位姿控制器以100 Hz以上的频率对位置和速度误差做出响应，其优化目标是使六自由度状态逼近给定参考值；而EMAT信号质量受升距波动、涡流分布变化和材料电磁属性共同调制，其退化模式无法单纯从位姿误差中推导。当前的主流做法是将控制回路和检测回路视为两个独立子系统——前者追求飞行稳定，后者被动接受信号波动——而缺少一个将两者耦合起来的中间层：一个能够联合解释位姿动力学、电磁-弹性耦合和视觉几何约束的多模态感知模型。

部分研究开始探索这一方向。基于视觉伺服的方法利用相机反馈估计探头-表面相对位姿，但其深度估计精度在近距离（<0.3 m）显著退化，恰与接触建立所需的亚厘米级定位精度构成矛盾。基于物理信息神经网络的方法通过在损失函数中嵌入动力学方程正则化项来约束模型输出，但现有工作几乎全部集中于单一系统的状态估计——如四旋翼的刚体动力学或流体力学——而尚未拓展到跨模态、跨时间尺度的异步观测融合问题。数据驱动的多模态融合方法（如基于Transformer的序列模型）在大规模语义任务中表现出色，但其注意力权重完全由特征相似度驱动，在面对物理测量噪声和传感器退化时可能产生不可解释甚至物理上矛盾的跨模态关联。

\subsubsection{本文的切入点}

上述分析勾勒出一条清晰的逻辑链：无人机检测已走完了从视觉遥感到物理接触的技术路径，但物理接触与有效信号耦合并非等价概念；EMAT因其对升距和姿态的多维敏感性，将这一矛盾推至极端；当前的解决方案要么聚焦于飞行控制（力控、全驱动），要么聚焦于信号处理（阈值、特征工程），而缺少一个将控制、感知和检测统一建模的系统性框架。本文正是在这一框架空缺处开展工作，提出以物理约束注意力为核心的多模态时序融合方法，将EMAT回波特征、位姿动态残差和RGBD视觉几何信息整合为具有物理可解释性的接触状态概率估计，为无人机EMAT自主检测填补感知-控制-检测协同闭环中的关键缺失环节。

\subsection{电磁超声厚度测量}

电磁超声换能器在非耦合条件下工作，其超声波直接在工件表面产生，相较于传统压电式超声换能器，能够避免工件表面粗糙和耦合层声速不均所引入的测量误差，且声速测量对换能器灵敏度要求较低，故EMAT的相关研究往往围绕厚度测量方向展开\citing{sun2025emat}。

涂君等人\citing{tu2021magnetic}提出了磁力可控电磁超声换能器来测量试件厚度，通过永磁体与电磁铁组合的方式控制磁场强度，进而实现磁场总体的增强和削弱。Chen W等人\citing{chen2024novel}提出了一种超声纵波共振方法，基于材料表面约束机制利用激光电磁声换能器来测量金属壁厚度，可在500°C高温和5.0 mm提离距离条件下实现误差为3.40\%的快速检测。Kubrusly等人\citing{kubrusly2023measurement}提出基于双电磁超声换能器的超声波最佳单向激发方法，通过双线圈干涉机制实现SH超声导波单向激发，测量误差低于0.8\%。郭伟等人\citing{guo2024metal}提出基于横波谐振频谱测厚的激光电磁超声技术，对三种金属材料在0.5-3.0 mm范围内的最大测厚偏差不超过6\%。苟浩瑞等人\citing{gou2024magnetic}提出多磁铁对称分布型EMAT，实现更小磁铁体积下更强的表面剩磁强度和更高的信噪比。Pei C等人\citing{pei2016modified}将单一永磁体结构改为双永磁体垂直方向反向排布实现磁场优化，使接收信号幅值提升一倍。

上述工作表明，EMAT测厚技术在静态实验条件下已达到较高精度。然而，这些方法均假设检测探头与被测表面处于稳定的相对位置关系，鲜有考虑探头动态运动对信号质量和测量精度的影响。在无人机搭载条件下，探头升距、法向角度和切向滑移的持续变化将直接影响涡流密度、磁耦合效率和回波可重复性，需要在动态物理约束下联合解释回波变化。

\subsection{无人机非侵入接触式检测}

近年来，多旋翼无人机在工业NDT领域的应用正经历从传统视觉巡检向接触式检测的深刻变革。接触式NDT要求无人机在与被测表面发生物理交互时，能够克服接触瞬间的冲击扰动，并维持稳定、精确的法向接触力。这一过程涉及复杂的运动学与动力学耦合，是当前空中机器人领域的研究热点与难点。纵观近年来的文献，该领域的研究主要沿着被动柔顺机构设计、主动力控与过驱动平台引入以及受限空间系统级框架集成三条主线展开。

在早期研究中，学者们主要侧重于通过机械结构设计来缓解接触碰撞带来的扰动。Gonz\'{a}lez-deSantos等人\citing{gonzalez2019payload}设计了一种具备碰撞缓解功能的柔性有效载荷，通过弹簧阻尼系统等被动柔顺机构有效吸收四旋翼无人机从自由飞行向接触状态过渡时的动能冲击，提升了系统在接触瞬间的鲁棒性。

随着对检测精度和接触力维持要求的不断提高，研究重心逐渐转向先进的主动控制策略。Kocer等人\citing{kocer2019inspection}针对天花板等上方结构检测，提出了一种基于非线性模型预测控制（NMPC）与非线性移动视界估计（NMHE）的框架，成功使四旋翼无人机在5.5 N的恒定接触力下稳定获取超声数据。

在硬件平台的进一步演进中，全驱动和过驱动无人机凭借独立的推力矢量控制，成为解决复杂面接触难题的另一核心路径。Ollero等人\citing{ollero2018aeroarms,trujillo2019novel}提出新型空中机械臂系统AeroX，通过在八旋翼无人机底部搭载多自由度机械臂实现检测探头与管壁的贴合。Rashad等人\citing{rashad2020towards}提出基于视觉的阻抗控制框架并应用于全驱动六旋翼无人机，利用全驱动平台位置与姿态解耦的运动学优势避免了传统欠驱动无人机在施加水平力时的姿态耦合问题。Watson等人\citing{watson2021deployment,watson2022dry}利用具有全矢量推力能力的Voliro多旋翼无人机，探讨了干耦合超声NDT的空中应用，实现了在倾斜和悬垂表面上的单点厚度测量和连续B扫描。Marcellini等人\citing{marcellini2024development}提出基于倾转旋翼的半自主NDT框架，通过双重阻抗控制架构动态吸收接触瞬间的冲击扰动。He\citing{he2025enhancing}提出了统一的空中操作框架，包含混合运动-力控制模块、接触感知轨迹规划器和以末端执行器为中心的控制接口。

在系统级应用阶段，Flyability公司推出了Elios 3 UT笼式无人机系统\citing{jamolli2025ndt,young2025mission}，将Cygnus双晶超声探头、智能机械臂与LiDAR传感器高度集成，能够在无需搭建脚手架的情况下安全深入受限空间进行贴壁超声检测。Veenstra等人\citing{veenstra2026autonomous}验证了商用四旋翼在非结构化环境中开展自主超声检测的可行性。

然而，上述研究主要聚焦于无人机从运动学角度实现“物理接触”，而鲜有探讨接触后“信号层面的可检测性”。真实的高空检测环境中，风场扰动和表面曲率突变会导致接触判定产生高频抖动，单纯的力学接触不等同于高信噪比的超声耦合。亟需将超声回波质量、位姿动态残差和视觉几何信息纳入统一的接触状态评价体系中。

\subsection{物理约束Transformer与多模态融合}

纯数据驱动模型具有“黑盒”特性，极易学习到数据中的伪关系，在面临分布外复杂扰动时容易输出违背物理常理的预测\citing{fakhro2025learning}。物理信息神经网络利用自动微分技术将常微分或偏微分方程作为物理残差项嵌入损失函数中，发挥了强大的正则化作用。Sanyal等人\citing{sanyal2023ramp}提出的RAMP-Net框架采用混合训练策略：物理损失负责捕获理想动力学方程以维持系统鲁棒性，数据损失自适应处理未建模的残余扰动。Gu等人\citing{gu2024physics}通过强制网络输出严格逼近6-DOF刚体运动学和牛顿-欧拉动力学方程，使PINN能有效过滤高频传感器噪声并填补稀疏未观测状态。

然而，基于MLP的PINN缺乏对时序上下文的显式建模。Zhao等人\citing{zhao2023pinnsformer}提出PINNsFormer框架，利用Transformer多头注意力机制精准捕捉全局时间依赖，并引入基于实数傅里叶变换的Wavelet激活函数以分解拟合复杂信号的高频振荡特征。

针对异构传感器采样频率不一导致的时间错位，Xie等人\citing{xie2026mpfusionnet}提出MPFusionNet，结合动态时间规整和三次样条插值，将低频视觉与高频雷达信号重采样至统一时序尺度，并通过基于PerceiverIO的交叉注意力机制实现多模态深层交互。Liu等人\citing{liu2025online}结合持续反向传播与PINN，提出在线状态估计框架，使无人机能够在保留离线学习的动力学规律前提下，通过在线小批量流式数据自适应外部未知气流的动态演变。

综合以上文献分析，现有的无人机接触式检测研究在三个维度上仍存在显著不足：（1）感知模态单一，接触状态判别主要依赖力学反馈或单一超声阈值，缺乏多模态联合表征；（2）物理约束缺失，纯数据驱动模型忽略了EMAT换能过程中电磁-弹性耦合的物理规律，泛化能力不足；（3）时序建模薄弱，接触概率的逐帧输出缺乏平滑与滞回机制，导致控制闭环中的状态闪烁。因此，构建多模态融合、物理约束嵌入与时序平滑约束一体化的接触状态感知框架，是提升无人机EMAT自主检测可靠性的关键方向。

\section{本文的主要贡献与创新}

针对上述研究不足，本文以多模态逐级物理约束为核心方法论，以无人机搭载EMAT自主探测为应用场景，开展了从物理建模、平台搭建、信号采集、特征融合到闭环决策的系统性研究。主要创新点与贡献如下：

\begin{enumerate}
\item \textbf{多模态接触状态感知框架}：首次将EMAT超声回波序列、无人机位姿动态残差及RGBD局部视觉ROI特征进行统一时序建模，将接触状态从瞬时二值标签提升为概率化的连续表达，克服了单模态阈值法的根本局限。
\item \textbf{物理约束注意力机制}：在Transformer自注意力计算中显式引入基于时间邻近和空间残差的物理约束矩阵，以对数偏置形式将无人机运动模型与接触物理先验嵌入跨模态特征交互过程，使注意力权重具有物理可解释性。
\item \textbf{多源不确定性传播与可靠性门控}：设计了模态可靠性门控权重，对视觉深度失真、位姿未收敛和超声信噪比下降等不同退化模式进行自适应证据分配，并通过双阈值持续时间约束状态机实现控制闭环的滞回平滑切换。
\item \textbf{完整实验验证体系}：搭建了集成LiDAR、RGBD相机、PX4飞控与EMAT探头的无人机多模态实验平台，开发了完整的ROS驱动与数据采集系统，并通过消融实验系统验证了物理约束与门控机制的增益。
\end{enumerate}

\section{本论文的结构安排}

本文共分为七章，章节结构安排如下：

第一章：绪论。阐述研究背景与意义，综述无人机NDT、EMAT测厚、接触式检测和物理约束Transformer四个方向的研究现状，明确本文的切入点和贡献。

第二章：电磁超声换能与接触检测理论基础。从电磁-弹性耦合机理出发，推导洛伦兹体力与弹性动力学方程的基本关系，建立接触边界与升距效应的回波统计表征模型。

第三章：无人机多模态实验平台构建。详细介绍无人机飞行平台、EMAT探头与ROS驱动系统、多传感器集成以及多模态同步数据采集系统的设计与实现。

第四章：多模态时序对齐与统一状态表示。分析多源异步采样问题，提出基于统一时间基准的加权插值对齐方法，设计各模态独立编码器并讨论传感器延迟补偿策略。

第五章：物理约束注意力机制与多模态融合。系统阐述物理约束矩阵的构造方法、物理约束注意力机制的设计和多源不确定性传播与可靠性判据框架。

第六章：实验验证与分析。描述实验设计与数据集构建方法，给出接触判别性能评估结果，并通过消融实验验证各模块功能，报告无人机闭环扫查的初步验证结果。

第七章：总结与展望。对全文工作进行总结，指出当前研究的局限性并展望未来研究方向。
